% 相対効果性の補助結果と普遍 cubic 表示の入力は分類論文
% \cite{QuadricIntrinsic2026} の対応する節から引用する。
% !TeX root = ../nc-quadric-surface-moduli.tex
% 2026-11-22 / v1.44（pointed Artin smooth bridgeとnoncubic chart次元を明確化）
\section{内在的な全モジュライの有限型代数性}
\label{sec:global-h2-smooth-atlas}

本節では、全ての noncubic 点で得た有限型の効果的近傍を用いる。
入力は実際の相対 $\mathfrak H_2$ 族であり、全幾何 fiber の分類だけではない。
これらの近傍から指定された pointed 同型を保つ smooth atlas を作り、
既存の noncubic 表示族の準コンパクト性を使って有限個の chart に縮約する。
有限 foundation スタックと内在的モジュライの位相を先に同一視する必要はない。

\subsection{noncubic 点を含む実際の smooth chart}

\begin{corollary}[効果的 noncubic 近傍の smooth atlas]
\label{cor:noncubic-effective-smooth-neighborhood}
$x\in\mathfrak H_2(k)$ を noncubic 点とする。
有限 foundation の標準表示空間 $S_0=\mathbb A^{64}$ の開部分 $S$ と、
$x$ の normal foundation を表す点 $s_x\in S(k)$ を取る。
定理~\ref{thm:noncubic-finite-type-effectivity} で構成した有限型 affine étale 近傍
\[
 a_x:(U_x,u_x)\longrightarrow(S,s_x)
\]
および実際の族 $q_x:U_x\to\mathfrak H_2$ を取る。
全軌道と canonical tails の比較により、この族には次の性質を持つ列 $(E_i)$ がある。
各 $E_i$ は平坦な noetherian compact 対象で、$E_0=\OO$ であり、
各四項窓は full strong exceptional collection、
全前向き derived Hom は次数零の有限 projective 加群である。
さらに $E_{i+4}=S^{-1}E_i[2]$ であり、最初の四項窓の End 代数は
頂点・積・任意基底変更を保って $a_x^*\Lambda$ と同定される。

このとき $q_x$ はスキームで表現可能な smooth 射である。
その fppf 像 $\mathfrak V_x\subset\mathfrak H_2$ は $x$ を含む開部分スタックであり、
有限型で滑らかな Artin スタックである。
\end{corollary}

\begin{proof}
定理~\ref{thm:atlas-from-effective-foundation} の仮定を確認する。
tails の row は次数成分ごとに基底上平坦であり、商関手は tensor と整合するので、
各 $E_i$ は平坦である。Noetherian 性と compact 性は実際の相対族の構成による。
四項の完全性・strong 性と前向き Hom の消滅は、五頂点 bar 比較により任意基底変更上で成立する。
full exceptional collection の Serre rotation は mutation による helix の延長と一致するため、
$E_{i+4}=S^{-1}E_i[2]$ は De Deken--Lowen の geometric $(4,3)$-helix 条件を与える。
最後の End 代数の同定も全軌道比較の一部である。
従って同定理を適用できる。

指定された同型に関する点を明記する。
small extension $A'\twoheadrightarrow A$ と
\[
 u_A\in U_x(A),\qquad c_{A'}\in\mathfrak H_2(A'),\qquad
 \alpha_A:q_x(u_A)\xrightarrow{\sim}c_{A'}|_A
\]
を取ると、$\alpha_A$ により四項窓とその全 Hom の基底を還元先へ移す。
De Deken--Lowen の持上げを基底 $A$ に直接適用し、
例外対象の $\Ext^1,\Ext^2$ の消滅と $A'^\times\to A^\times$ の全射性により、
$E'_0=\OO_{c_{A'}}$ とその指定された pointing を保つ。
五つの Hom 空間の基底も指定還元を延長して選ぶと、実際の合成から $S(A')$ の点が得られる。
$a_x$ の全点での étale 性により、これは $u_A$ を保つ $u_{A'}\in U_x(A')$ へ持ち上がる。

両方の四項 tilting 窓の End 代数をこの基底で同定する。
得られる標準 dg Morita 同値の還元は、対象だけでなく全ての射上で $D(\alpha_A)$ と一致する。
補題~\ref{lem:artin-two-heart-comparison}、または反 Serre 再構成の全忠実性を使えば、
これは $\alpha_A$ 自体を延長する pointed exact 同型となる。
module-object 基底変更により、その同型は全有限型 test 上の相対系にも延長する。
この議論は根や被覆を同型の標識へ加えず、指定した pointed 同型をそのまま保つ。

これらの指定された Artin 同型の持上げと、
$\Phi:\mathfrak H_2\to\mathfrak P$ の全忠実性は
補題~\ref{lem:pointed-artin-lifting-smooth-bridge} の全仮定を満たす。
従って $q_x$ はスキームで表現可能な smooth 射で、
その開像は全ての pointed 同型を relation に含めた
smooth groupoid により表される。
この推論は有限型スキーム上の Tag 02HY の Artin 試験に帰着し、
現時点で $\mathfrak H_2$ 全体が Artin スタックであるとは仮定しない。
\end{proof}

\subsection{有限個の chart による全体の表示}

\begin{theorem}[全 $\mathfrak H_2$ の有限型代数性と滑らかさ]
\label{thm:global-h2-smooth-finite-type}
局所有限型 $k$-スキームの site 上の内在的スタック $\mathfrak H_2$ は、
$k$ 上有限型で滑らかな純次元三の Artin スタックの制限により表される。
その対角は affine かつ有限表示である。
有限型の滑らかなスキーム $U$ と representable smooth 全射
\[
 q:U\longrightarrow\mathfrak H_2
\]
を取れ、全ての pointed exact 同値を用いる relation scheme
\[
 R=U\times_{\mathfrak H_2}U
   \simeq U\times_{\mathfrak P}U
\]
に対して $\mathfrak H_2\simeq[U/R]$ である。
ここで選ばれた foundation、Serre 根、平方同型、AS4 被覆を保存することは、射の条件に加えない。
\end{theorem}

\begin{proof}
まず有限個の noncubic 近傍で足りることを示す。
定理~\ref{thm:noncubic-relative-effectivity} による実際の射を、
表示スタックの atlas へ引き戻すと
\[
 p:P=Q^{\mathrm{nd}}\times Q^{\mathrm{nd}}
   \longrightarrow\mathfrak N^{\mathrm{pres}}
   \longrightarrow\mathfrak H_2
\]
を得る。$P$ は有限型の準コンパクトな $k$-スキームである。
命題~\ref{prop:noncubic-presentation-stack} により、全 noncubic $k$-対象はその像に入る。
各 noncubic 点 $x$ に前系の開部分 $\mathfrak V_x$ を取る。
引戻し $p^{-1}(\mathfrak V_x)$ は $P$ の開部分スキームであり、これらは $P$ の全閉点を覆う。
有限型 $k$-スキームは Jacobson なので、全体も覆う。
従って準コンパクト性により有限個
\[
 p^{-1}(\mathfrak V_{x_1}),\ldots,p^{-1}(\mathfrak V_{x_m})
\]
で $P$ を覆える。この段階で使ったのは $P\to\mathfrak H_2$ の存在と
各 $\mathfrak V_x$ の内在的な開性だけであり、$\mathfrak N^{\mathrm{pres}}\to\mathfrak H_2$
の full faithfulness や、任意の相対 noncubic 族の正規形は使わない。

定理~\ref{thm:cubic-smooth-open} の有限型 smooth atlas $U_{\mathrm{cub}}$ と合わせ、
\[
 U=U_{\mathrm{cub}}\amalg U_{x_1}\amalg\cdots\amalg U_{x_m}
\]
と置く。有限和なので $U$ は有限型で滑らかである。
各成分の射は representable smooth なので $q:U\to\mathfrak H_2$ も同様であり、
その開像は全ての $k$-対象を含む。

全ての相対族を覆うことも確認する。
有限型 affine $T\to\mathfrak H_2$ に対して、$U\times_{\mathfrak H_2}T\to T$ の像は開である。
$T$ の各閉点の剰余体は $k$ であり、その fiber は上の分類と有限被覆によって像に入る。
従ってこの開像は全閉点を含み、Jacobson 性により $T$ 全体である。
ゆえにこの射は smooth 全射であり、fppf 局所に section を持つ。
任意の局所有限型 $T$ については有限型 affine 開被覆へ制限すればよい。
これで $q$ が元の site 上の全対象を fppf に覆うことが示された。

相対反 Serre 認識の全忠実性から、relation は表示した $U\times_{\mathfrak P}U$ と等しい。
$\mathfrak P$ の対角が affine 有限表示なので、$R$ は有限型スキームである。
両射影 $R\rightrightarrows U$ は $q$ の基底変更として smooth であり、
\cite[Tags 04X0, 04TK]{StacksSmoothAtlas} により $[U/R]$ は通常の全スキーム上の Artin スタックとなる。
その元の site への制限は $\mathfrak H_2$ と同値である。
有限型 atlas $U$ によって有限型性が従い、$U$ の滑らかさによって $k$ 上の滑らかさが従う。
対角の性質は命題~\ref{prop:relative-h2-algebraicity-reduction} による。

最後に、系~\ref{cor:cubic-formal-h2-abelian} と系~\ref{cor:noncubic-all-artin-h2} により、
各 $k$-点の形式変形 groupoid は pointed Abel 変形と同値である。
その接空間は $\operatorname{HH}^2(\Lambda)$、無限小自己同型空間は
$\operatorname{HH}^1(\Lambda)$ であり、
命題~\ref{prop:foundation-hochschild} によって両次元の差は三である。
滑らかな Artin スタックの局所次元はこの差に等しいから、全閉点で三となる。
有限型 smooth atlas の各成分は閉点を持ち、次元は局所一定なので、全体は純次元三である。
\end{proof}

\begin{corollary}[noncubic 有効 chart の相対次元]
\label{cor:noncubic-atlas-relative-dimension}
上の $q_x:U_x\to\mathfrak H_2$ は、像上で相対次元
$64-3=61$ の smooth 射である。
その relation scheme
$R_x=U_x\times_{\mathfrak H_2}U_x$ は
$k$ 上 smooth な純次元 $64+61=125$ のスキームである。
\end{corollary}
\begin{proof}
$U_x$ は $S\subset\mathbb A^{64}$ の有限型 étale 近傍なので、
その各点における次元は64である。
定理~\ref{thm:global-h2-smooth-finite-type} により
$\mathfrak H_2$ は純次元3でsmoothである。
smoothなスタック射の相対次元公式より $q_x$ は相対次元61。
$R_x\to U_x$ は $q_x$ の基底変更であるため、
相対次元61の smooth 射で、総次元は125となる。
これはsmooth atlasのfiber次元の整合性を表すものであり、
非cubic点のモジュライに新たな変形方向を加えたものではない。
\end{proof}

\subsection{普遍 cubic 族の降下と内在的な既約性}

\begin{lemma}[根を標識にしない普遍非線型 cubic 族]
\label{lem:universal-nonlinear-cubic-h2-family}
$S_{\mathrm e}=\mathscr G\setminus\mathscr G_{\mathrm{lin}}$ を、
第~\ref{sec:relative-picard-roots}節の非線型 geometric 五つ組の開部分とする。
普遍 $\mathbb Z$-代数 $D$ と $C_{ij}=D_{-j,-i}$ から、実際の相対族
\[
 q_{\mathrm e}:S_{\mathrm e}\longrightarrow\mathfrak V_{\mathrm{cub}}
\]
が得られる。この射は representable smooth であり、その開像は既約である。
開像の補集合の幾何点は、存在するとすれば可換二次曲面の pointed 圏 $Q$ だけである。
\end{lemma}

\begin{proof}
系~\ref{cor:finite-type-root-chart} の étale 被覆
$\coprod_\alpha T_\alpha\to S_{\mathrm e}$ を取る。
各 $T_\alpha$ には実際の二周期があり、定理~\ref{thm:geometric-cubic-relative-effectivity} により
$q_\alpha:T_\alpha\to\mathfrak H_2$ が構成されている。
その構成における偶数 restriction は、任意の基底変更上で
\[
 \operatorname{QGr}B_\alpha\simeq\operatorname{QGr}C_{T_\alpha}
\]
という構造対象を保つ exact 同値を与える。
右辺の $C$ は根を選ぶ前の普遍代数であり、二周期の選択に依存しない。
従って $T_\alpha\times_{S_{\mathrm e}}T_\beta$ 上では、両族は同じ
$\operatorname{QGr}C$ を通じて pointed exact 同値となる。
三重の重なりの cocycle は、この同じ圏を通る比較の結合性による。
coherent 対象と全有限型 module-object 基底変更にも整合するので、
$\mathfrak H_2$ の fppf descent により $q_{\mathrm e}$ が得られる。
根と AS4 被覆は局所的 witness であり、それら自体を重なりで同定する必要はない。

各 $q_\alpha$ は系~\ref{cor:quintuple-effective-smooth-chart} により smooth である。
smoothness は始域の étale 被覆で判定できるから $q_{\mathrm e}$ も smooth である。
$S_{\mathrm e}$ は射影空間の非空開部分で既約なので、その連続像は既約であり、
smooth 性によって像は開である。
全非線型 cubic 表示が像に入ることは、定理~\ref{thm:cubic-smooth-open} の逆構成による。
線型 geometric cubic 表示は全て同じ pointed $Q$ を与えるので、最後の主張が従う。
\end{proof}

\begin{proposition}[全内在的モジュライの既約性]
\label{prop:global-h2-irreducible}
$\mathfrak V_{\mathrm{cub}}$ は既約であり、$\mathfrak H_2$ において稠密である。
従って $\mathfrak H_2$ は既約かつ連結である。
\end{proposition}

\begin{proof}
前補題の既約な開像を $\mathfrak V_{\mathrm e}$ と書く。
残り得る点 $Q$ をその閉包に入れる。
系~\ref{cor:quadric-effective-smooth-neighborhood} の $U_Q$ で、指定点 $u_Q$ を含む
既約成分 $U_Q^0$ に制限する。
$U_Q$ は smooth なので、その既約成分は開である。
$U_Q^0\to S\subset S_0=\mathbb A^{64}$ は étale であり、像は非空開部分である。

線型 geometric cubic 表示の四項 foundation は、可換二次曲面の
\[
 (\OO,\OO(1,0),\OO(1,1),\OO(2,1))
\]
および ruling を交換した列から得られる。
従ってその $S_0$ 上の基底付き表示は、有限個の $\mathcal G/K$ 軌道に含まれる。
$\dim(\mathcal G/K)=61<64$ なので、この有限和の閉包は $S_0$ の真の閉部分である。
さらに命題~\ref{prop:cubic-foundation-open} の非空な geometric cubic 開部分と交差する。
この二つの稠密開条件を同時に課すと、$U_Q^0$ の稠密開部分の foundation は非線型 cubic であり、
その $q_Q$ による像は $\mathfrak V_{\mathrm e}$ に入る。
実際の射 $q_Q$ の連続性から、$Q$ は $\mathfrak V_{\mathrm e}$ の閉包に入る。
前補題の像が最初から $Q$ を含む場合にも同じ結論である。
従って $\mathfrak V_{\mathrm{cub}}$ は、この一つの既約な開部分の閉包であり、既約である。

任意の noncubic $k$-点 $x$ に対し、
定理~\ref{thm:noncubic-finite-type-effectivity} の $U_x$ のうち指定点を含む既約成分を取る。
これも $S_0$ の非空開部分への étale 射を持つ。
命題~\ref{prop:cubic-foundation-open} の geometric cubic foundation は
$S_0$ の非空開部分なので、$U_x$ の稠密開部分は cubic である。
その実際の族 $q_x$ は全幾何 fiber の foundation と heart を保つため、
この稠密開部分は $\mathfrak V_{\mathrm{cub}}$ へ送られる。
従って $x$ は $\mathfrak V_{\mathrm{cub}}$ の閉包に入る。

以上で全 $k$-点がこの閉包に入った。
もし補集合が非空なら、有限型 smooth atlas への逆像は非空開部分であり、
$k$-点を持つ。この像も $k$-点なので矛盾する。
従って $\mathfrak V_{\mathrm{cub}}$ は $\mathfrak H_2$ において稠密である。
既約部分空間の閉包は既約なので結論を得る。
この証明で用いた連続写像は全て実際の相対 $\mathfrak H_2$ 族の射であり、
finite foundation の幾何点像の位相を内在位相と同一視する操作はない。
\end{proof}

\subsection{noncubic 境界の滑らかさと余次元}

\begin{lemma}[二つの射影行列式の横断性]
\label{lem:lagrangian-boundary-transverse}
$W\subset\operatorname{LG}(4,U\oplus U^*)$ を
命題~\ref{prop:lagrangian-geometric-recognition-open} の認識開部分とし、
tautological 四平面束の二射影を
\[
 M:\mathscr L\longrightarrow U\otimes\OO_W,
 \qquad N:\mathscr L\longrightarrow U^*\otimes\OO_W
\]
と書く。
\[
 D=V(\det M,\det N)\subset W
\]
は $k$ 上 smooth な純次元八のスキームであり、$D\hookrightarrow W$ は
余次元二の正則閉埋込みである。
$D$ 上の交点線束を
\[
 \ell=\mathscr L\cap U=\ker N,\qquad
 m=\mathscr L\cap U^*=\ker M
\]
とすると、その法束は
\[
 N_{D/W}\simeq(m^*)^{\otimes2}\oplus(\ell^*)^{\otimes2}
 \tag{NC-boundary-normal}\label{eq:noncubic-boundary-normal}
\]
と同定される。
\end{lemma}

\begin{proof}
認識開条件により、$D$ の全幾何点で $M,N$ はともに階数三である。
三次小行列式を可逆化して行列を掃き出すと、射影は $\operatorname{diag}(I_3,a)$ と書ける。
$D$ 上ではその四次行列式が実際に零なので $a=0$ である。
従って非被約構造を先に除かなくても kernel と cokernel は線束であり、
$\ell,m$ は $\mathscr L$ の線部分束となる。
二 summand の交わりが零なので $\ell\oplus m\subset\mathscr L$ は階数二の部分束である。

幾何点 $L\in D$ において、標準交代形式から
\[
 T_L\operatorname{LG}\simeq\operatorname{Sym}^2L^*
\]
を得る。
具体的には一次変形 $\varphi:L\to(U\oplus U^*)/L$ に対し、
持上げ $\widetilde\varphi$ を用いた
$q(v,w)=\omega(v,\widetilde\varphi(w))$ が対称形式となる。
$q$ は持上げの選択に依存しない。

$\operatorname{im}M=m^\perp\subset U$ なので、
$\operatorname{coker}M\simeq m^*$ である。
行列式を微分すると、階数三の射 $M$ に対するその一次変分は
$\ker M=m$ から $\operatorname{coker}M=m^*$ への誘導射で計算される。
局所的には、これは $\operatorname{diag}(I_3,0)$ の右下成分の微分にほかならない。
従って $\det M$ の微分は、符号を含む単元倍を除いて
\[
 \operatorname{Sym}^2L^*\longrightarrow(m^*)^{\otimes2},
 \qquad q\longmapsto q|_{m\times m}
\]
である。同様に $\det N$ の微分は $q|_{\ell\times\ell}$ である。
二線は独立なので、この二つの restriction を合わせた射は全射である。
例えば $\ell,m$ の生成元を $L$ の基底の最初の二元とすれば、
二つの写像は対称行列の二つの独立な対角係数を取る。

行列式線束を局所自明化すれば、$D$ は滑らかな十次元スキーム $W$ 上の
二つの関数で定義され、その微分は全幾何点で独立である。
Jacobian 判定により $D$ は smooth で余次元二、両関数は正則列となる。
特にこの証明は $D$ の reducedness も同時に与える。

最後に行列式線束を intrinsic に読むと、kernel と cokernel の完全列から
\[
 (\det U\otimes\det\mathscr L^*)|_D
 \simeq\det(\operatorname{coker}M)\otimes\det(\ker M)^*
 \simeq(m^*)^{\otimes2}
\]
であり、$N$ 側は $(\ell^*)^{\otimes2}$ である。
上の微分の全射とその kernel の接束への同定から
\eqref{eq:noncubic-boundary-normal} を得る。
\end{proof}

\begin{theorem}[内在的 noncubic 境界は滑らかな余次元二の閉部分]
\label{thm:noncubic-boundary-smooth-codim-two}
定理~\ref{thm:global-h2-smooth-finite-type} の有限型 Artin スタックにおいて、
\[
 \mathfrak N=
 \bigl(\mathfrak H_2\setminus\mathfrak V_{\mathrm{cub}}\bigr)_{\mathrm{red}}
\]
を reduced induced closed substack とする。
このとき $\mathfrak N$ は有限型で滑らかな既約 Artin スタックであり、純次元一である。
閉埋込み $\mathfrak N\hookrightarrow\mathfrak H_2$ は余次元二の正則閉埋込みである。
具体的に、各 noncubic 点を含む適当な smooth scheme chart では、
そのイデアルは実際の二つの射影行列式 $(\det M,\det N)$ で生成される。
\end{theorem}

\begin{proof}
\emph{境界の定義。}
$\mathfrak V_{\mathrm{cub}}$ は内在的な開部分であり、$\mathfrak H_2$ は有限型なので、
その閉補集合には一意な reduced induced closed substack がある
\cite[Tag 0509, Lemma 100.10.1]{StacksReducedBoundary}。
これを $\mathfrak N$ とする。

\emph{quotient chart での余次元二。}
指定した noncubic 点に対し、良い根に対応する Lagrangian 表示を取る。
命題~\ref{prop:noncubic-lagrangian-etale-slice} の étale locus と認識開条件に制限し、
$G^+$-不変開 $W$ と
\[
 \mathfrak L=[W/G^+]\longrightarrow\mathfrak F
\]
を得る。ここで $\mathfrak F$ は該当する finite foundation スタックの開部分であり、
この射は representable étale である。
$G^+$ は二つの summand を保つ tensor 作用と反傾作用で作用するので、
射影の行列式は局所的に単元倍される。
従ってイデアル $(\det M,\det N)$ と $D$ は $G^+$-不変である。
前補題と smooth atlas $W\to[W/G^+]$ により
\[
 \mathfrak D=[D/G^+]\hookrightarrow\mathfrak L
\]
は余次元二の正則閉埋込みであり、$\mathfrak D$ は smooth である。
その次元は $8-\dim G^+=8-7=1$ である。
ここで正則性の降下は Koszul-regular immersion の fpqc locality を用い、
局所 Noetherian な表示上で通常の regular immersion と同値であることを使う
\cite[Tags 0694, 063L]{StacksReducedBoundary}。

\emph{両方の smooth 性を保つ共通 chart。}
定理~\ref{thm:noncubic-finite-type-effectivity} と
系~\ref{cor:noncubic-effective-smooth-neighborhood} の
$q_0:U_0\to\mathfrak H_2$ を取る。
$U_0\to S$ は étale、$S\to\mathfrak F$ は foundation の smooth atlas なので、
$U_0\to\mathfrak F$ は smooth である。
そこで
\[
 V=U_0\times_{\mathfrak F}\mathfrak L
\]
を取り、指定点を含む étale scheme chart $U\to V$ に移る。
このとき $V\to U_0$ は étale、$V\to\mathfrak L$ は smooth であるから、
\[
 q:U\longrightarrow\mathfrak H_2,
 \qquad a:U\longrightarrow\mathfrak L
\]
はともに smooth である。
必要ならさらに étale に細分して $G^+$-torsor を自明化し、
tautological 束の基底を選べば、実際の行列 $M,N$ を $U$ 上で取れる。
この細分後も二つの smooth 性は保たれる。

\[
 D_U=U\times_{\mathfrak L}\mathfrak D
\]
と置く。これは $U$ 上で $V(\det M,\det N)$ と書ける閉部分スキームである。
$a$ は smooth なので、前段の正則閉埋込みを flat に引き戻している。
従って $D_U\hookrightarrow U$ は余次元二の正則閉埋込みとなり、
$D_U\to\mathfrak D$ が smooth であることから $D_U$ 自身も $k$ 上 smooth である。
この議論は、torsor の切断から得る lift $U\to W$ の smooth 性や flat 性を仮定しない。

\emph{内在的境界との scheme-theoretic な同定。}
一方、$\mathfrak N_U=U\times_{\mathfrak H_2}\mathfrak N$ は $U$ の閉部分スキームである。
射 $\mathfrak N_U\to\mathfrak N$ は $q$ の基底変更として smooth であり、
$\mathfrak N$ は reduced なので $\mathfrak N_U$ も reduced である。
ここでは reducedness が smooth topology に局所的であることを使う
\cite[Tag 034E]{StacksReducedBoundary}。

認識開条件の下では、両射影が階数三の幾何点は noncubic normal foundation を与え、
一方が階数四なら一方の ruling foundation が geometric cubic となる。
実際の相対族はその fiber heart と構造対象を保つため、
系~\ref{cor:lagrangian-open-closed-recognition} と体上の分類により
\[
 |\mathfrak N_U|=|D_U|
\]
である。
両者は同じスキーム $U$ の reduced closed subscheme なので、
reduced induced 構造の一意性から
\[
 \mathfrak N_U=D_U
\]
が closed subscheme として成立する。
幾何点の一致から非被約な二つの部分スキームを無条件に同一視しているのではなく、
前補題と smooth pullback によって双方の reducedness を確認した後の同定である。

全 noncubic 点でこの共通 chart を取れるので、これらの $D_U$ は $\mathfrak N$ を smooth に覆う。
従って $\mathfrak N$ は $k$ 上 smooth であり、
$\mathfrak N\hookrightarrow\mathfrak H_2$ は smooth 局所で余次元二の正則閉埋込みである。
$\mathfrak H_2$ は純次元三なので $\mathfrak N$ は純次元一となる。

\emph{既約性。}
定理~\ref{thm:noncubic-relative-effectivity} の実際の族を
\[
 P=Q^{\mathrm{nd}}\times Q^{\mathrm{nd}}\longrightarrow\mathfrak H_2
\]
へ引き戻す。$P$ は $\mathbb A^6$ の非空開部分であり、reduced かつ既約である。
その全幾何 fiber は noncubic なので、この射は $\mathfrak N$ を経由する。
この因数分解には、reduced な始域から閉部分への set-theoretic な包含が
scheme-theoretic な因数分解を与えることを使う
\cite[Tag 050B, Lemma 100.10.3]{StacksReducedBoundary}。
体上の noncubic 分類を各代数閉拡大体上で適用すると、$P\to\mathfrak N$ は全幾何点を覆う。
従って $|\mathfrak N|$ は既約空間 $|P|$ の連続像であり、既約である。
この結論に、$\mathfrak N^{\mathrm{pres}}\to\mathfrak N$ の相対全忠実性や同値は用いない。
\end{proof}

\begin{remark}[非被約 test 上で境界に属する条件]
$\mathfrak N$ は reduced induced closed substack である。
非被約な test $T\to\mathfrak H_2$ が $\mathfrak N$ を経由することは、
全幾何 fiber が noncubic であることだけからは従わない。
上の smooth chart では、前者は $\det M=\det N=0$ が
$\OO_T$ の元として実際に成立することを要求する。
行列式が非零の冪零元となる横断的な Artin 変形は、閉 fiber が noncubic でも境界には属さない。
また本定理は reduced boundary の正則埋込みを証明するものであり、
それ以外の候補 functor に非被約構造を無条件に与える主張ではない。
\end{remark}

\subsection{framed 実現空間と ambient 内の非開性}

\begin{corollary}[framed 実現空間のスキーム性]
\label{cor:global-xas-smooth-scheme}
\[
 X_{\mathrm{as}}=X_3\times_{\mathfrak P}\mathfrak H_2
\]
は有限型で滑らかな既約・純次元八十三のスキームにより表される。
$X_{\mathrm{as}}\to X_3$ は有限型の monomorphism で、separated かつ unramified、quasi-finite である。
\end{corollary}

\begin{proof}
全スキーム上では前定理の $[U/R]$ を用いる。
$R=U\times_{\mathfrak P}U$ は全ての同型を記録するため、
$[U/R]\to\mathfrak P$ は全スキーム上でも全忠実である。
実際、対象を $U$ へ fppf 局所に持ち上げると Isom はまさに $R$ であり、
この一致は降下する。従って $X_{\mathrm{as}}\to X_3$ は monomorphism である。

命題~\ref{prop:relative-h2-algebraicity-reduction} により、$X_{\mathrm{as}}$ は有限型代数空間である。
また $X_{\mathrm{as}}\to\mathfrak H_2$ は $\operatorname{PGL}_9$-torsor なので smooth で相対次元八十である。
従って $X_{\mathrm{as}}$ は $k$ 上 smooth で純次元 $3+80=83$ である。
命題~\ref{prop:global-h2-irreducible} により $\mathfrak H_2$ は既約で、
この torsor は開全射かつ幾何 fiber が既約な $\operatorname{PGL}_9$ である。
従って $X_{\mathrm{as}}$ も既約である。
実際、二つの非空開部分の像は底の既約性により交わり、
その fiber 内の二つの非空開部分は fiber の既約性により交わる。

$X_{\mathrm{as}}$ と $X_3$ はともに有限型 $k$-空間なので、両者の間の射は有限型である。
局所有限型 monomorphism は separated、unramified、locally quasi-finite かつスキームで表現可能である
\cite[Tag 05W6, Lemma 76.14.8]{StacksRepresentableMonomorphism}。
特に $X_3$ がスキームであるため $X_{\mathrm{as}}$ 自身もスキームである。
同じ結論は separated locally quasi-finite な代数空間からスキームへの射に関する
\cite[Tag 03XX, Proposition 67.50.2]{StacksRepresentableMonomorphism} からも従う。
有限型性は準コンパクト性を含むので、射は quasi-finite である。
\end{proof}

\begin{proposition}[canonical 実現像は ambient 全体で開ではない]
\label{prop:global-h2-not-ambient-open}
$\Phi:\mathfrak H_2\to\mathfrak P=[X_3/\operatorname{PGL}_9]$ は、
noncubic の二つの根対が一致する点で open immersion ではない。
従って $X_{\mathrm{as}}\to X_3$ も、その点の任意の frame の近傍で open immersion ではない。
\end{proposition}

\begin{proof}
一致点 $x$ において、系~\ref{cor:noncubic-all-artin-h2} は
\[
 \dim T_x\mathfrak H_2=4,\qquad
 \dim\operatorname{HH}^1(\Lambda)=1
\]
を与える。ここで $T_x$ は閉 fiber の同定を保つ一次変形の同型類の空間である。
一方、系~\ref{cor:artin-polarization-retraction} の分裂完全列は
\[
 0\longrightarrow\operatorname{HH}^1(\Lambda)
 \longrightarrow T_x\mathfrak P
 \longrightarrow\operatorname{HH}^2(\Lambda)\longrightarrow0
\]
を与え、$\dim T_x\mathfrak P=5$ である。
$\Phi$ は canonical 反 Serre 環を取る射なので、その微分の像は
この分裂における $\operatorname{HH}^2$ の四次元の像である。
従って一次変形の射は全射ではなく、$\Phi$ はこの点で open immersion にならない。
開 immersion なら形式変形 groupoid を同一にすることに反するからである。

framed 空間については、その任意の frame $J$ で
\[
 \dim T_JX_{\mathrm{as}}=83,\qquad
 \dim T_JX_3=5+(80-1)=84
\]
となる。後者の $80-1$ は $\operatorname{PGL}_9$ の軌道接空間の次元である。
open immersion ならこの二つの接空間は同型になるので、framed 射についても結論を得る。
\end{proof}

\begin{remark}[表現可能性から自動的には出ない比較]
上で構成したのは内在的な族のスタックそのものの有限型 smooth atlas と、
framed 部分層のスキームによる表現である。
これによって $X_{\mathrm{as}}$ を $X_3$ の開部分と同一視することはできず、
実際に前命題が一致根対でその同一視を否定する。
また有限 foundation の ruling 商 $\mathfrak Q_{\mathrm{fnd}}$ との相対同値や、
その表示からの内在位相の比較は、この atlas の構成だけでは結論しない。
全スキーム上の $[U/R]$ は元の有限型 site 上の $\mathfrak H_2$ を延長するが、
任意の Noetherian 基底上で素朴に同じ Abel 公理を課した関手との一致は別の連続性の問題である。
\end{remark}

\subsection{偏極を変えた関係式の開近傍}
\label{subsec:twisted-polarization-neighborhood}

$V=k^9$、$G=\operatorname{PGL}(V)$ と置く。以下の $J$ は
$X_{\mathrm{as}}$ の点と、それが $V^{\otimes2}$ 内に定める関係部分束を表す。
系~\ref{cor:global-xas-smooth-scheme} の monomorphism により、後者を固定すれば
前者の選択は高々一つである。有限型のスキーム
\[
 D=\{(J,g)\in X_{\mathrm{as}}\times G:
                      (g\otimes g)J=J\}
\]
を取る。等式は Grassmannian の対角の逆像であり、スキームとして課す。
射影行列を線型行列に持ち上げても、スカラーは部分束を変えない。
$D$ は $X_{\mathrm{as}}$ 上で pointed 自己同値を表す群スキームである
（命題~\ref{prop:relative-h2-stack-recognition}）。特に平坦性は要求しない。

\begin{theorem}[次数移動を記録する開埋込み]
\label{thm:polarization-inertia-open}
次の $G$-同変射は開埋込みである。
\[
 \theta:D\longrightarrow X_3,\qquad
 (J,g)\longmapsto (g^{-1}\otimes 1)J .
\]
始域で $G$ は関係式の枠を変え、$g$ には共役で作用する。
像を $W$ とすると $[W/G]$ は $\HHstack$ の慣性スタックに同型である。
$X_{\mathrm{as}}\to X_3$ は $W$ の中の恒等切断 $g=1$ に一致し、
有限表示の局所閉埋込みである。
\end{theorem}
\begin{proof}
\emph{写像と逆構成の一意性。}
枠と次数の線束を局所自明化する。$R=T(V)/(J)$ の次数付き自己同型として
$g$ を持ち上げると、積 $a*b=g^{\deg b}(a)b$ の右 Zhang twist は同じ
次数成分を持ち、関係部分束は $(g^{-1}\otimes1)J$ である。
実際、長さ $n$ の語の積は順に $g^{n-1},\ldots,g,1$ を掛けてから元の積を
取るので、二次関係の全挿入もこの同型で対応する。
スカラーの取り方によらず $X_3$ への射が得られる。
右加群の明示同値は命題~\ref{prop:antiserre-polarization-family} の
式~\eqref{eq:polarization-right-zhang} の証明と同じで、任意の係数変更と可換する
\cite[Theorem 3.1]{Zhang1996}。

圏の言葉では、$(\cC,\OO)$ の pointed 自己同値 $u$ が反 Serre 環に
誘導する射影行列が $g$ である。$t=S^{-1}[2]$ と $u$ の可換同型は
Serre 双対の自然性から得られる。$u\OO\simeq\OO$ を使うと、
$F=ut$ の軌道対象 $F^n\OO$ は $t^n\OO$ と同定される。
右規約の積 $F^{\deg b}(a)\circ b$ をこの同定で移したものが上記の積である。
負の軌道対象も一致し、従って AZ の二条件は $t$ から $F$ へ移る。
全軌道環の tails は、構造対象と表示された移動を含めて
$(\cC,\OO,F)$ を再現する。
逆にこの triple から $t$ を内在的に回収すると $u=Ft^{-1}$ が決まる。
$u\OO\simeq\OO$ の選択はスカラーだけ異なり、同じ pointed 射の同型類を与える。
また、移動を比較する自然同型の違いは
補題~\ref{lem:relative-h2-flat-center} によりスカラーで、
次数 $n$ にはその $n$ 乗だけ作用する。従って、もとの次数一の
\emph{射影枠}も回収できる。具体的には
$\Hom(\OO,F\OO)\simeq\Hom(\OO,t\OO)$ に枠を移す。
線束 twist のある局所自明化の重なりでもこの射影同定は変わらない。
以上で同じ framed twist から二つの異なる $(J,g)$ は得られない。
特に任意の代数閉拡大体上で単射である。

\emph{非恒等 twist の Artin 近傍。}
$k$ 点 $(J_0,g_0)$ において、剰余体 $k$ の局所 Artin 環 $A$ 上の
任意の framed 関係式変形 $J'_A$ を取り、
$B_A=T(V_A)/(J'_A)$、$F_A=(1)$ と書く。
閉 fiber $B_0$ は Koszul である。実際、命題~\ref{prop:all-h2-antiserre-koszul}
の環の単純加群の線型自由解消を twist すれば、同じ生成次数の自由解消を得る。
左右の次数付き noetherian 性も twist で保たれる。
$X_3$ の次数三の平坦性と補題~\ref{lem:relative-koszul-degree-three} から
$B_A$ は全次数自由である。補題~\ref{lem:artin-noetherian-lift} と
命題~\ref{prop:artin-relations-flat-abelian} の torsion 商の証明を適用する。
同証明が使うのは全次数平坦性、閉 fiber の noetherian 性、
および次数一生成性であり、これらはいま確認した。
よって $(\qgr B_A,\pi B_A)$ は元の $(\cC_0,\OO_0)$ の flat pointed Abel 変形である。
系~\ref{cor:cubic-formal-h2-abelian}、\ref{cor:noncubic-all-artin-h2}
によって、これは Artin 基底上で実際の相対 $\HHstack$ 族になり、
指定した還元同型も保たれる。

この族の反 Serre を $t_A$ とする。$u_A=F_A t_A^{-1}$ は閉 fiber で
$u_0$ となる。$u_A\OO_A$ と $\OO_A$ は同じ例外対象の平坦持上げである。
$\Ext^1(\OO_0,\OO_0)=0$ と小拡大による帰納法で閉 fiber の指定同型を持ち上げられる。
得られた pointed 自己同値は加群対象への基底変更で系全体に延びる。
従って $F_A=u_A t_A$ の軌道対象は $t_A$ のものと同定され、
その全軌道 Hom は自由で基底変更と可換する。
自然射 $B_A\to\bigoplus_{n\ge0}\Hom(\OO_A,F_A^n\OO_A)$ は
閉 fiber で上の twist の同型である。各次数が同じ階数の自由加群なので
Nakayama によって全次数で同型となる。
ゆえに次数一の枠を移した canonical 関係式と $u_A$ の行列が
$(J_A,g_A)\in D(A)$ を与え、$\theta(J_A,g_A)=J'_A$ である。

この逆構成は Artin 基底変更と整合する。実際、平坦な $B_A$ の
torsion 商の係数還元は、商の右随伴を用いる
命題~\ref{prop:artin-relations-flat-abelian} の議論を
$A\to A/I$ に対して適用すれば、$\operatorname{QGr}(B_A/IB_A)$ である。
$F_A$ と自由 row の還元もこれに従う。圏と $t_A$ の比較は相対認識の
全忠実性による。逆に $D(A)$ から始めれば、twist の pointed tails 同値と
射影枠の回収により元の点に戻る。特に小拡大で還元先の $D$ の点を
あらかじめ指定していても、その点を保つ一意の持上げとなる。

\emph{開埋込みと切断。}
両スキームは代数閉体 $k$ 上有限型なので、局所 Artin 小拡大による
smoothness 試験 \cite[Tag 02HY]{StacksSmoothAtlas} の有限型閉点の剰余体は $k$ である。
上の持上げにより $\theta$ は smooth で、双対数での一意性により
閉点で相対接空間は零である。従って全域でエタールである。
幾何点上の単射性と \cite[Tag 025G]{StacksBoundaryComparison}
により開埋込みを得る。$[D/G]$ は $\HHstack=[X_{\mathrm{as}}/G]$
の慣性であり、$D$ の恒等切断は $G$ の単位元の逆像だから閉である。
\end{proof}

\begin{corollary}[任意の関係式族への適用]
\label{cor:x3-foundation-neighborhood}
局所有限型基底 $S$ と任意の $j:S\to X_3$ に対し、$S_W=j^{-1}W$
は、その幾何 fiber が構造対象を保つ自己同値による反 Serre 環の
Zhang twist となる点の開集合である。$S_W$ の fppf 局所で階数二十四の
normal foundation $\Lambda$ と、表示された次数移動を実現する
可逆 perfect $\Lambda$-双加群を構成できる。canonical 実現の部分は
$S_W$ 上で $g=1$ を課して得られ、任意の基底変更に整合する。
\end{corollary}
\begin{proof}
開埋込みの逆を $S_W$ 上で取ると、唯一の canonical 族とその自己同値が得られる。
定理~\ref{thm:global-h2-smooth-finite-type} の cubic・noncubic の
smooth atlas を引き戻し、その相対 tilting 窓の和 $T$ と $\Lambda=\End(T)$ を取る。
twist の tails 同値で移した $F=ut$ に対して
$P=\mathbf R\Hom(T,F(T))$ と置く。左作用は $F$ が $T$ の射に及ぼす作用、
右作用は前合成である。標準の導来増強で $F$ は余極限を保ち、
自由生成対象上の値とこの作用から関手は $-\otimes_\Lambda^{\mathbf L}P$
となる。$F^{-1}$ が逆双加群を与える。$P$ は基底上 perfect であり、
$\Lambda^e$ の相対 smooth 性により双加群としても perfect である
（$\Lambda^e$ の有限対角解消に $P$ を tensor する）。
相対 compact Hom の基底変更からこの表示は引戻しと整合する。
次数の正規化による基底線束は局所自明化してこの記述を適用する。
恒等切断を引き戻すと所要の閉条件を得る。行列代表を取れば、
非対角成分と対角成分相互の差をすべて零にする方程式である。
非被約な基底ではこれらの式が基底環で消えることを要求する。
\end{proof}

\begin{example}[$W$ の外の三次表示]
九生成元 $x_i$ に対して、二文字 $x_ix_j$ は $1\le j\le d_i$ のときだけ残し、
その他を零とする。ただし $(d_1,\ldots,d_9)=(0,3,4,0,3,3,4,4,4)$ とする。
長さ二の語は $25$ 個である。$x_i$ で始まる長さ三の語数は
$\sum_{j\le d_i}d_j$ で、$i=1,4$ で零、残りの七箇所で七である。
同じ漸化式により長さ四の語数は $98$ となる。
従って $X_3$ の点を与えるが、$W$ 上の四次成分の階数は $81$ である。
この例では $x_2,x_3$ 以外を零とする商が二生成自由代数になるので、
右 noetherian 性も成り立たない。
\end{example}
